%! TEX root = main.tex
\begin{feedback}[\cref{ex:la-7}]
Your conclusions for all three maps are correct. In part (a), the matrix
is also correct: its columns are the images of $e_1$ and $e_2$. The
linearity calculation starts with $T_1(au+bv)$ only implicitly, so label
the two sides or introduce the expression you are expanding. The same
issue occurs in part (c), where the first displayed line names $T$ rather
than $T_3$ and then writes $T_1$ and $T_2$ on the right. Keep the map
names consistent. There are also unmatched right parentheses after
$au_1-2au_2$ and $bv_1-2bv_2$; remove these rather than leaving them as
algebraic noise.

Part (b) is a good counterexample. To make the calculation fully
transparent, identify the left side as $T_2(2x)$ and the right side as
$2T_2(x)$, rather than only showing the resulting coordinates. Part (c)
should not be described as similar to $T_2$ ``with a smaller rank'': $T_2$
is nonlinear, so its rank has not been defined in this lesson, and the
reason $T_3$ is linear is the coordinate calculation you give.
\end{feedback}

\begin{feedback}[\cref{ex:la-8}]
Your recovery of $T(e_1)$ and $T(e_2)$ from the two given values is
correct, and the general formula for $T(x,y)$ follows correctly from
those images. There is one sign error in the matrix: you found
$T(e_2)=(1/2,-1/2,1/2)^T$, but its third entry becomes $+1/2$ in your
matrix instead of $-1/2$. Thus the displayed matrix does not represent
the formula immediately above it.

When you revise, use each recovered standard-basis image as a column and
check your matrix against both of the originally supplied values. The
phrase ``where the basis vectors are in the image'' is also imprecise:
you have found their images, not located the basis vectors inside
$\operatorname{im}T$.
\end{feedback}

\begin{feedback}[\cref{ex:la-9}]
The row reduction, pivot columns, equations for the kernel, and image
basis are correct. In particular, you correctly took the image basis from
the original matrix. Your two kernel vectors do belong to $\ker T$, and
they are linearly independent: for example, their third and fourth
coordinates form the matrix
\[
  \begin{pmatrix}1&2\\-1&-6\end{pmatrix},
\]
whose determinant is nonzero. They also have the right number for a
two-dimensional kernel.

The missing step is the derivation of these vectors from the free
variables. Set $x_3=s$ and $x_4=t$, solve the two displayed equations for
$x_1$ and $x_2$, and show that the resulting general vector is a linear
combination of your two claimed vectors. Without that calculation, the
sentence about two degrees of freedom explains the expected dimension,
but does not prove that this particular pair spans the entire kernel.
The notation $\ker A$ is acceptable for a matrix map, but $\ker T$ matches
the exercise's notation.
\end{feedback}

\begin{feedback}[\cref{ex:la-10}]
The conclusion is correct, and your use of linearity followed by
injectivity contains the central idea. However, the proof does not yet
handle the definition of linear independence for the whole list. A
dependent list need not contain a vector $v_c$ that is a combination of
just two other vectors $v_a$ and $v_b$. It might instead satisfy a
relation involving three or more vectors, so your three-vector equation
only treats a special case.

Start instead with an arbitrary relation
$\sum_{i=1}^k c_iT(v_i)=0$, where $c_1,\ldots,c_k\in\R$. Use linearity
and injectivity to obtain a relation among $v_1,\ldots,v_k$, and then
state what their independence forces about every $c_i$. You can prove
the result directly; contradiction is not necessary.
\end{feedback}

\begin{feedback}[\cref{ex:la-11}]
You have correctly proved that injectivity implies surjectivity in the
equal finite-dimensional setting. Writing
$\operatorname{span}(\operatorname{im}T)\ne W$ is redundant because
$\operatorname{im}T$ is already a subspace; the direct statement is
$\operatorname{im}T\ne W$. The step from $\dim\ker T>0$ to
noninjectivity needs the reason that a nontrivial kernel contains a
nonzero vector $v$, while $T(v)=0=T(0)$. Listing distinct vectors
$v_1,\ldots,v_k$ does not itself establish that there are at least two.

More importantly, this is an ``if and only if'' exercise, and the
converse is missing. You must also assume that $T$ is surjective and use
$\dim\operatorname{im}T=n$ together with rank--nullity to determine the
kernel. In both directions, state explicitly where the equal dimensions
are used, and use the kernel criterion for injectivity from the lesson.
\end{feedback}

\begin{feedback}[\cref{ex:la-12}]
Your first observation about a nonzero kernel vector is on the right
track, but make the collision explicit: if $v\ne0$ and $T(v)=0$, then
$T(x+v)=T(x)$ for every input $x$. Thus $x$ and $x+v$ are distinct inputs
with the same output, which explains why the output cannot determine the
input uniquely. Merely saying that the input ``can't be reconstructed''
could be read as a claim about practical computation rather than this
mathematical nonuniqueness.

The rank discussion needs correction. The scalar coordinates of $Ax$ are
always real numbers; saying that they are ``not linearly independent''
when $\operatorname{rank}A<m$ is not the relevant statement. Instead,
describe the image as a subspace of $\R^m$ of dimension
$\operatorname{rank}A$: at most that many independent output degrees of
freedom can vary, even though the output has $m$ scalar entries. Also,
rank--nullity constrains rank by the number of input columns; it does not
by itself justify your claim involving $m$.

The final composition argument is not correct as written. Nonnegative
rank and nullity do not imply that neither can exceed $m$, and
``more expressive ($\operatorname{rank}A>m$) outputs'' is impossible
for a map $\R^n\to\R^m$: the rank is always at most $m$ because
$\operatorname{im}A\subseteq\R^m$. The rank is also at most $n$, by
rank--nullity. Replace this part with an argument about the composition
itself: two matrices $A$ and $B$ compose to another single matrix, and
the image of the composition is contained in the image of the outer map.
Explain why this bounds, rather than increases, the available output
degrees of freedom.
\end{feedback}
